% !TEX root = ../main.tex
\chapter{Giriş}
\label{chap:giris}

Bu bölüm problemi ve tezin katkılarını tanıtır. Kısaltmalar ilk kullanımda
açılır: \ac{API}, ikinci kullanımda kısa hâli görünür: \ac{API}.

Tezin geri kalanı şöyle düzenlenmiştir: \cref{chap:literatur} ilgili
çalışmaları, \cref{chap:yontem} önerilen yöntemi, \cref{chap:bulgular}
deneysel bulguları ve \cref{chap:sonuc} sonuçları sunar.

\section{Motivasyon}
\label{sec:motivasyon}

Kaynak gösterimi \textcite{knuth1984texbook} ve parantez içinde
\parencite{lamport1994latex} şeklinde yapılır.

\section{Katkılar}
\label{sec:katkilar}

\begin{itemize}
  \item Birinci katkı.
  \item İkinci katkı.
\end{itemize}
